\section*{Reading guide}
\addcontentsline{toc}{section}{Reading guide}

The guide has four parts. Part I gives a stand-alone route from data to
results. Parts II and III form the detailed guide: Part II develops the
economic argument, and Part III turns that argument into data, software, and
applications. Part IV collects technical material that is easier to consult
separately.

\subsection*{Document structure}

\textbf{Part \ref{part:quickstart}: Quick start.}
Section \ref{ch:quickstart} is a stand-alone operating procedure for a new
application. It distinguishes the economic-geography and urban-commuting
models, states the input accounting conditions, shows the required file
schemas and configuration choices, and gives the commands for validation,
analysis, and decomposition. Begin there when the immediate question is whether
an existing network dataset can be used with the package.

\textbf{Part \ref{part:theory}: Detailed guide---theory and welfare.}
Section \ref{ch:question} fixes the policy experiment, the three reported
welfare measures, and the boundary between model-implied benefits and project
appraisal. Section \ref{ch:equilibrium} then defines the economic-geography and
urban models together with their shared transport equations. Section
\ref{ch:welfare} uses those definitions to establish the traffic result,
derive the extended welfare elasticity, and decompose the difference between
the two approaches.

\textbf{Part \ref{part:implementation}: Detailed guide---implementation and
applications.}
Section \ref{ch:implementation} states the data contract, follows the software
calculation, and gives the numerical and reporting checks. Section
\ref{ch:worked} applies the complete sequence to the public multimodal grid.
Section \ref{ch:applications} distinguishes the audited U.S. freight results,
whose replication requires external data, from the urban and Seattle
extensions. Section \ref{ch:conclusion} collects the interpretation and
reporting requirements for a result.

\textbf{Part \ref{part:appendices}: Reference appendices.}
The appendices are reference material rather than another stage of the main
argument. Appendix \ref{app:crosswalk} maps notation to code; Appendix
\ref{app:algorithm} records the computational sequence; Appendix
\ref{app:config} gives the configuration and schema reference; and Appendix
\ref{app:verification} states the numerical tests and tolerances. Appendices
\ref{app:additional-examples} and \ref{app:provenance} summarize the other
package examples and document the sources behind the guide.

\subsection*{Suggested reading paths}

\begin{center}
\small
\begin{tabularx}{\textwidth}{L{0.25\textwidth}X}
\toprule
Reader's task & Suggested path \\
\midrule
Adapt a new dataset &
Section \ref{ch:quickstart}; then Sections \ref{ch:implementation} and
\ref{ch:worked}; keep Appendices \ref{app:config} and
\ref{app:verification} open while preparing the inputs. \\
Assess the economics &
Sections \ref{ch:question}--\ref{ch:welfare}; then Appendix
\ref{app:crosswalk} for the equation-to-code map. \\
Build an urban or public-transit application &
The model comparison in Section \ref{ch:quickstart}; the urban model and shared
transport equations in Section \ref{ch:equilibrium}; then the urban material
in Section \ref{ch:applications}. \\
Extend the package &
Sections \ref{ch:equilibrium}--\ref{ch:implementation}; then Appendices
\ref{app:crosswalk}--\ref{app:verification}. \\
Audit or reproduce a result &
Sections \ref{ch:implementation}, \ref{ch:worked}, and
\ref{ch:applications}; then Appendices \ref{app:verification} and
\ref{app:provenance}. \\
\bottomrule
\end{tabularx}
\end{center}

\subsection*{Conventions used throughout}

The economic-geography and urban models share routing, modal choice,
congestion, and primitive-cost pass-through, but they have different spatial
states and welfare projections. Public transit is therefore a set of modes in
the common transport block, not a third spatial model. The guide uses three
labels consistently. The \emph{Traditional approach} uses traffic on the
improved edge--mode.
The \emph{full realized-cost effect} changes the cost users face after
congestion. The \emph{Extended approach} is the full primitive-cost effect of
the infrastructure-side improvement. Model-ready flows are an input
requirement. Vehicle counts, passenger counts, schedules, or tonnage require a
documented conversion and balancing step before they enter the package.
